\documentclass[a4paper,11pt]{article}

% --- CJK support (xelatex) ---
\usepackage{fontspec}
\usepackage{xeCJK}
\setCJKmainfont{Yu Mincho}
\setCJKsansfont{Yu Gothic}
\setCJKmonofont{Yu Gothic}

% --- Packages ---
\usepackage{amsmath,amssymb}
\usepackage{geometry}
\usepackage{graphicx}
\usepackage{booktabs}
\usepackage{hyperref}
\usepackage{fancyhdr}
\usepackage{xcolor}
\usepackage{enumitem}
\usepackage{tabularx}
\usepackage{array}

% Column type L: auto-wrapping left-aligned column for tabularx
\newcolumntype{L}{>{\raggedright\arraybackslash}X}

\geometry{margin=2.0cm}
\setlength{\headheight}{14pt}
\hypersetup{colorlinks=true, linkcolor=blue, urlcolor=blue}

\pagestyle{fancy}
\fancyhf{}
\rhead{SaunaFlow}
\lhead{支配方程式ドキュメント}
\cfoot{\thepage}

\title{\textbf{SaunaFlow 支配方程式ドキュメント}}
\author{SaunaFlow Project}
\date{\today}

\begin{document}
\maketitle
\tableofcontents
\newpage

\section{SaunaFlow 支配方程式ドキュメント}

本ドキュメントでは、SaunaFlow で使用する \textbf{2 つのソルバー} それぞれの支配方程式、
離散化手法、ゾーン構成、数値解法を整理する。

\bigskip\hrule\bigskip

\bigskip\hrule\bigskip

\subsection{1. ソルバー概要}

{\small
\noindent
\begin{tabularx}{\textwidth}{|L|L|L|}
\hline
\textbf{項目} & \textbf{正確版 (OpenFOAM)} & \textbf{簡易版 (simple\_solver.py)} \\
\hline
次元 & 3D & 0D (2ゾーン集中定数) \\
\hline
方程式 & N-S + エネルギー + 乱流 & ゾーン質量保存 + エネルギー保存 + プルーム相関式 \\
\hline
乱流モデル & SST k-omega & なし (プルーム相関式に内包) \\
\hline
圧力-速度連成 & PIMPLE 法 (非定常) & なし (速度場を解かない) \\
\hline
時間積分 & backward (2次精度陰解法) & 前進 Euler (擬似時間進行) \\
\hline
用途 & 本計算・検証 & 可視化 UI・パラメータスタディ \\
\hline
\end{tabularx}
}


\bigskip\hrule\bigskip

\subsection{2. 正確バージョン — OpenFOAM (buoyantPimpleFoam)}

\subsubsection{2.1 連続の式 (質量保存)}

$
\frac{\partial \rho}{\partial t} + \nabla \cdot (\rho \mathbf{U}) = 0
$

\begin{itemize} \item $\rho$: 密度 [kg/m³]（完全ガス状態方程式で計算） \end{itemize}
\begin{itemize} \item $\mathbf{U}$: 速度ベクトル [m/s] \end{itemize}
\begin{itemize} \item 非定常計算 (buoyantPimpleFoam) のため時間微分項を保持 \end{itemize}

\subsubsection{2.2 運動量保存 (Navier-Stokes)}

$
\frac{\partial (\rho \mathbf{U})}{\partial t}
+ \nabla \cdot (\rho \mathbf{U} \otimes \mathbf{U})
= -\nabla p_{rgh}
+ \nabla \cdot \left[ \mu_{\text{eff}} \left( \nabla \mathbf{U} + (\nabla \mathbf{U})^T - \frac{2}{3} (\nabla \cdot \mathbf{U}) \mathbf{I} \right) \right]
+ \rho \mathbf{g}
$

ここで:

\begin{itemize} \item $p_{rgh} = p - \rho \mathbf{g} \cdot \mathbf{r}$: 静水圧を除いた修正圧力 [Pa] \end{itemize}
\begin{itemize} \item $\mu_{\text{eff}} = \mu + \mu_t$: 実効粘性 [Pa·s] \end{itemize}
\begin{itemize} \item $\mathbf{g} = (0, -9.81, 0)$ m/s²: 重力加速度 \end{itemize}
\begin{itemize} \item \textbf{浮力駆動流}: 密度差 → 圧力勾配 → プルーム上昇 \end{itemize}

離散化 (\texttt{fvSchemes}):

\begin{verbatim}
div(phi,U):  bounded Gauss linearUpwind grad(U)   ← 2次風上 (数値拡散低減)
grad(U):     cellLimited Gauss linear 1            ← セル制限付き線形
ddt:         backward                              ← 2次精度陰的時間積分
\end{verbatim}

\subsubsection{2.3 エネルギー保存}

$
\frac{\partial (\rho h)}{\partial t}
+ \nabla \cdot (\rho \mathbf{U} h) = \nabla \cdot (\alpha_{\text{eff}} \nabla T) + q_{\text{source}}
$

\begin{itemize} \item $h = c_p T$: 比エンタルピー (sensibleEnthalpy) [J/kg] \end{itemize}
\begin{itemize} \item $\alpha_{\text{eff}} = \frac{\mu}{Pr} + \frac{\mu_t}{Pr_t}$: 実効熱拡散率 [W/(m·K)] \end{itemize}
\begin{itemize} \item $q_{\text{source}}$: ヒーター熱流束 [W/m²]（境界条件として付与） \end{itemize}
\begin{itemize} \item $c_p = 1005$ J/(kg·K), $Pr = 0.7$ \end{itemize}

離散化:

\begin{verbatim}
div(phi,T):  bounded Gauss linearUpwind default    ← 2次風上
laplacian:   Gauss linear corrected                ← 2次中心差分 + 非直交補正
\end{verbatim}

\subsubsection{2.4 乱流モデル (SST k-omega)}

Menter の SST (Shear Stress Transport) k-omega モデルを採用。
標準 k-epsilon に比べ、壁面近傍の低レイノルズ数領域と
自然対流における熱伝達率の予測精度が優れる。

\textbf{乱流運動エネルギー k:}

$
\frac{\partial (\rho k)}{\partial t}
+ \nabla \cdot (\rho \mathbf{U} k)
= \nabla \cdot \left[ (\mu + \sigma_k \mu_t) \nabla k \right]
+ P_k + G_b - \beta^* \rho k \omega
$

ここで $G_b$ は浮力生成項:

$
G_b = -\frac{\mu_t}{\rho \, Pr_t} \left( \mathbf{g} \cdot \nabla \rho \right)
$

\begin{itemize}
\item $Pr_t = 0.85$ (乱流プラントル数)
\item $\mathbf{g} = (0, -9.81, 0)$ m/s$^2$ (重力加速度ベクトル)
\item 安定成層 (上部高温・下部低温): $G_b < 0$ → 乱流を抑制し、温度層の鮮明化に寄与
\item 不安定成層 (ヒーター近傍): $G_b > 0$ → 乱流を促進し、対流混合を強化
\item 実装: 修正乱流モデルではなく fvOptions の codedSource により k 方程式に注入
\end{itemize}

\textbf{比散逸率 omega:}

$
\frac{\partial (\rho \omega)}{\partial t}
+ \nabla \cdot (\rho \mathbf{U} \omega)
= \nabla \cdot \left[ (\mu + \sigma_\omega \mu_t) \nabla \omega \right]
+ \frac{\gamma}{\nu_t} P_k - \beta \rho \omega^2
+ 2(1 - F_1) \frac{\rho \sigma_{\omega 2}}{\omega} \nabla k \cdot \nabla \omega
$

\textbf{$\omega$ 方程式における浮力項の扱い:}
現在の実装では $G_b$ は $k$ 方程式にのみ注入し、$\omega$ 方程式への浮力項（$C_3 \gamma G_b / \nu_t$ 等）は
未追加である。これは OpenFOAM 標準の \texttt{kOmegaSST} が $\omega$ に対する
浮力ソースを持たない設計に準じている。
強い成層下では $\omega$ 方程式への浮力散逸項が界面の乱流抑制精度に影響するため、
Phase 2 以降で \texttt{buoyancyDissipation} の有効化を検討する。

\textbf{乱流粘性:}

$
\mu_t = \frac{\rho a_1 k}{\max(a_1 \omega, \; S F_2)}
$

ここで $F_1$, $F_2$ はブレンディング関数、$S$ はひずみ速度テンソルの大きさ。

\textbf{SST k-omega の利点 (サウナ解析において):}

\begin{itemize} \item 壁面近傍: k-omega モデルが自動的に適用され、壁面関数への依存が軽減 \end{itemize}
\begin{itemize} \item 主流域: k-epsilon 相当の挙動にブレンドされ安定 \end{itemize}
\begin{itemize} \item 自然対流の壁面熱伝達率を k-epsilon より正確に予測 \end{itemize}

\textbf{壁面境界条件:}

{\small
\noindent
\begin{tabularx}{\textwidth}{|L|L|}
\hline
\textbf{変数} & \textbf{壁面 BC} \\
\hline
k & kqRWallFunction \\
\hline
omega & omegaWallFunction \\
\hline
nut & nutkWallFunction \\
\hline
alphat & compressible::alphatWallFunction (Prt = 0.85) \\
\hline
\end{tabularx}
}


離散化:

\begin{verbatim}
div(phi,k):      bounded Gauss upwind    ← 1次風上 (安定性重視)
div(phi,omega):  bounded Gauss upwind    ← 1次風上
\end{verbatim}

\subsubsection{2.5 状態方程式}

$
\rho = \frac{p M_w}{R T}
$

\begin{itemize} \item $M_w = 28.96$ g/mol (空気の分子量) \end{itemize}
\begin{itemize} \item $R = 8314$ J/(kmol·K) \end{itemize}
\begin{itemize} \item 完全ガス (perfectGas) 仮定 \end{itemize}
\begin{itemize} \item 熱物性モデル: \texttt{heRhoThermo} + \texttt{pureMixture} + \texttt{hConst} + \texttt{const} transport \end{itemize}

\textbf{浮力の発生メカニズム:}
ヒーター近傍で温度上昇 → 密度低下 → 浮力 → 上昇プルーム形成

\subsubsection{2.6 時間平均場}

非定常計算のため、統計量として時間平均場を出力する。
\texttt{fieldAverage} 関数オブジェクトにより、指定開始時刻以降の
$\bar{T}$, $\bar{U}$, $\overline{T'^2}$, $\overline{U'^2}$ を計算。

\begin{verbatim}
averaging_start: end_time * 0.5  (デフォルト: 後半50%を平均化)
\end{verbatim}

\bigskip\hrule\bigskip

\subsection{3. 簡易バージョン — 2-Zone プルームモデル}

\subsubsection{3.1 モデル概要}

建築環境工学で標準的な \textbf{2層ゾーンモデル} を採用。
Morton-Taylor-Turner (MTT) のエントレインメント理論と
Zukoski のプルーム相関式に基づく。

\begin{verbatim}
 ┌─────────────────────────────────┐  ← ceiling
 │                                 │
 │   上層 (Upper hot layer)         │  温度: T_upper
 │   密度: rho_upper = rho_0 * T_wall / T_upper
 │                                 │
 ├ ─ ─ ─ ─ ─ ─ ─ ─ ─ ─ ─ ─ ─ ─ ─ ┤  ← 界面高さ z_int
 │                                 │
 │   下層 (Lower cool layer)        │  温度: T_lower
 │   密度: rho_lower = rho_0 * T_wall / T_lower
 │          ┌───┐                  │
 │          │ H │ ← ヒーター         │
 │          │   │   (プルーム発生源)   │
 └──────────┴───┴──────────────────┘  ← floor
              ↑
           プルーム上昇
         (エントレインメント)
\end{verbatim}

\textbf{1D 拡散モデルとの根本的な違い:}

\begin{itemize} \item 対流を「大きな熱拡散」で近似する代わりに、プルームの質量・エネルギー輸送を直接モデル化 \end{itemize}
\begin{itemize} \item 質量保存則を遵守（プルーム上昇量 = 壁面沿い下降量） \end{itemize}
\begin{itemize} \item 物理的根拠のある相関式（Zukoski のプルーム相関）を使用 \end{itemize}

\subsubsection{3.2 プルームモデル (Zukoski 相関式)}

ヒーター中心から高さ $z$ の位置におけるプルーム質量流量:

$
\dot{m}_p(z_{\text{eff}}) = 0.071 \, Q_c^{1/3} \, z_{\text{eff}}^{5/3} + 0.0018 \, Q_c
$

ここで $z_{\text{eff}}$ は Heskestad の仮想原点補正を適用した有効高さ:

$
z_0 = -1.02 \, D + 0.083 \, Q_c^{2/5}
$

$
z_{\text{eff}} = \max(0.01, \; z - z_0)
$

\begin{itemize} \item $D$: ヒーター特性直径 [m] (YAML の \texttt{heater.width}) \end{itemize}
\begin{itemize} \item $z_0$: 仮想原点オフセット [m] (小型ヒーターでは通常負 → $z_{\text{eff}} > z$) \end{itemize}

Zukoski 相関は点源に対して導出されたため、有限径 $D$ のヒーターには
Heskestad (1984) の仮想原点補正が必要。これにより、実効プルーム高さが
幾何学的高さより大きくなり、エントレインメント量が適切に増加する。

ここで:

\begin{itemize} \item $Q_c$: 対流熱出力 [kW] ($Q_c = f_{\text{conv}} \times Q_{\text{heater}}$, $f_{\text{conv}} = 0.7$) \end{itemize}
\begin{itemize} \item $z$: ヒーター中心からの鉛直距離 [m] \end{itemize}
\begin{itemize} \item $\dot{m}_p$: プルーム質量流量 [kg/s] \end{itemize}

プルーム温度（エネルギー保存から）:

$
T_{\text{plume}} = T_{\text{lower}} + \frac{Q_c}{\dot{m}_p \, c_p}
$

\textbf{物理的意味:} プルームが上昇するにつれ周囲空気をエントレイン（巻き込み）し、
質量流量は増加するが温度は低下する。これは MTT 理論の本質的な特徴。

\subsubsection{3.3 支配方程式}

\textbf{状態変数:} $T_{\text{upper}}$ (上層温度), $T_{\text{lower}}$ (下層温度), $z_{\text{int}}$ (界面高さ)

\paragraph{上層エネルギー保存}

$
\rho_{\text{upper}} \, c_p \, V_{\text{upper}} \, \frac{dT_{\text{upper}}}{dt}
= \underbrace{\dot{m}_p \, c_p \, (T_{\text{plume}} - T_{\text{upper}})}_{\text{プルームからの熱入力}}
- \underbrace{h_{\text{wall}} \, A_{\text{wall,upper}} \, (T_{\text{upper}} - T_{\text{wall}})}_{\text{壁面熱損失}}
+ \underbrace{Q_{\text{rad}} \, F_{\text{upper}}}_{\text{ヒーター輻射 (fixed wall時)}}
$

ここで $F_{\text{upper}} = F_{\text{ceiling}} + F_{\text{upper\_walls}}$ はヒーターから上層壁面への形態係数の合計。
\texttt{model: lumped} の場合、輻射は壁面集中定数モデル（付録C）を経由して
壁内面温度 $T_{\text{wall,inner}}$ を加熱し、間接的に上層空気を温めるため、この直接項は 0 となる。

ここで:

\begin{itemize} \item $V_{\text{upper}} = A_{\text{floor}} \times (H - z_{\text{int}})$: 上層体積 [m³] \end{itemize}
\begin{itemize} \item $A_{\text{wall,upper}} = P \times (H - z_{\text{int}}) + A_{\text{floor}}$: 上層が接する壁面積 (側壁 + 天井) [m²] \end{itemize}
\begin{itemize} \item $P = 2(W + D)$: 水平断面の周長 [m] \end{itemize}

\paragraph{下層エネルギー保存}

$
\rho_{\text{lower}} \, c_p \, V_{\text{lower}} \, \frac{dT_{\text{lower}}}{dt}
= \underbrace{Q_{\text{rad}} \times 0.3}_{\text{輻射受熱}}
+ \underbrace{k_{\text{int}} \, A_{\text{floor}} \, \frac{T_{\text{upper}} - T_{\text{lower}}}{0.1 H}}_{\text{界面伝導}}
- \underbrace{h_{\text{wall}} \, A_{\text{wall,lower}} \, (T_{\text{lower}} - T_{\text{wall}})}_{\text{壁面熱損失}}
$

ここで:

\begin{itemize} \item $Q_{\text{rad}} = (1 - f_{\text{conv}}) \times Q_{\text{heater}}$: ヒーターからの輻射熱 [W] \end{itemize}
\begin{itemize} \item $k_{\text{int}} = 0.5$ W/(m·K): 界面の実効熱伝導率 \end{itemize}

\paragraph{界面質量保存}

$
A_{\text{floor}} \, \frac{dz_{\text{int}}}{dt}
= -\frac{\dot{m}_p}{\rho_{\text{upper}}} + \dot{V}_{\text{return}} - \dot{V}_{\text{steam}} + \dot{V}_{\text{mix}}
$

ここで:

\begin{itemize} \item $\dot{m}_p / \rho_{\text{upper}}$: プルームが上層に供給する体積流量（界面を押し下げる） \end{itemize}
\begin{itemize} \item $\dot{V}_{\text{return}}$: 壁面沿い下降流による体積流量（界面を押し上げる） \end{itemize}
\begin{itemize} \item $\dot{V}_{\text{steam}}$: ロウリュ蒸気による上層体積膨張（界面を押し下げる） \end{itemize}
\begin{itemize} \item $\dot{V}_{\text{mix}}$: アウフグース双方向質量交換による正味体積効果 \end{itemize}

$
\dot{V}_{\text{return}} = \frac{h_{\text{wall}} \, A_{\text{wall,upper}} \, (T_{\text{upper}} - T_{\text{wall}})}{\rho_{\text{upper}} \, c_p \, \max(T_{\text{upper}} - T_{\text{lower}}, \, 1)}
$

アウフグース質量交換項: $\beta_{\text{aug}}$ [kg/s] の双方向混合では、
上層空気が下層へ、下層空気が上層へ同時に移動する。
密度差により等質量の体積が異なるため、正味の体積変化が生じる:

$
\dot{V}_{\text{mix}} = \beta_{\text{aug}} \left( \frac{1}{\rho_{\text{lower}}} - \frac{1}{\rho_{\text{upper}}} \right)
$

$\rho_{\text{upper}} < \rho_{\text{lower}}$（高温空気は低密度）のとき $\dot{V}_{\text{mix}} < 0$ となり、
界面が下降する。これはタオルが高温空気を強制的に下方へ押す効果を表す。

定常状態ではこれらがバランスし、界面高さが安定する。

\paragraph{密度の温度依存性}

$
\rho = \rho_0 \frac{T_{\text{wall}}}{T}
$

\begin{itemize} \item $\rho_0 = 1.1$ kg/m³ (基準密度, \textasciitilde{}300K) \end{itemize}

\subsubsection{3.4 温度プロファイルの構築}

2-Zone の集中定数 ($T_{\text{upper}}$, $T_{\text{lower}}$) から
鉛直方向の連続プロファイルをシグモイド関数で補間:

$
T(y) = T_{\text{lower}} + \frac{T_{\text{upper}} - T_{\text{lower}}}{1 + \exp\left(-3 \cdot \frac{y - z_{\text{int}}}{\delta / 2}\right)}
$

ここで $\delta = 0.15 H$ は遷移層の厚さ。

\subsubsection{3.5 境界条件と制約}

{\small
\noindent
\begin{tabularx}{\textwidth}{|L|L|}
\hline
\textbf{境界/制約} & \textbf{条件} \\
\hline
界面高さ & $0.05 H \le z_{\text{int}} \le 0.95 H$ \\
\hline
上層温度 & $T_{\text{wall}} \le T_{\text{upper}} \le T_{\text{wall}} + 200$ K \\
\hline
下層温度 & $T_{\text{wall}} - 1 \le T_{\text{lower}} \le T_{\text{upper}}$ \\
\hline
\end{tabularx}
}


\subsubsection{3.6 参考文献}

\begin{itemize} \item Morton, Taylor \& Turner (1956), "Turbulent Gravitational Convection from Maintained and Instantaneous Sources", Proc. R. Soc. A 234:1-23 \end{itemize}
\begin{itemize} \item Zukoski (1978), "Development of a Stratified Ceiling Layer in the Early Stages of a Closed-Room Fire", NBS-GCR-78-150 \end{itemize}
\begin{itemize} \item Cooper (1982), "A Mathematical Model for Estimating Available Safe Egress Time in Fires", NBSIR 82-2612 \end{itemize}
\begin{itemize} \item Heskestad (1984), "Engineering Relations for Fire Plumes", Fire Safety Journal 7(1):25-32 \end{itemize}

\bigskip\hrule\bigskip

\subsection{4. ゾーン構成と熱収支}

\subsubsection{4.1 正確版 (OpenFOAM) のゾーン構成}

\begin{verbatim}
              fixedValue T_wall
         ┌──────────────────────────┐
         │        ceiling            │
         │                          │
         │     ┌──────┐             │
  heater │     │ 内部 │             │ fixedValue
  _wall  │     │ 領域 │             │ T_wall
 (熱流束)│     │      │             │ (opposite_wall)
         │     │      │             │
         │     └──────┘             │
         │  heater_wall_surround    │
         │   (fixedValue T_wall)    │
         └──────────────────────────┘
                  floor
              fixedValue T_wall

  front/back: fixedValue T_wall
\end{verbatim}

\textbf{各パッチの境界条件:}

{\small
\noindent
\begin{tabularx}{\textwidth}{|l|L|L|L|}
\hline
\textbf{パッチ名} & \textbf{温度 BC} & \textbf{速度 BC} & \textbf{圧力 BC} \\
\hline
heater\_wall & externalWallHeatFluxTemperature (flux) & noSlip & fixedFluxPressure \\
\hline
heater\_wall\_surround & fixedValue $T_{\text{wall}}$ & noSlip & fixedFluxPressure \\
\hline
floor & fixedValue $T_{\text{wall}}$ & noSlip & fixedFluxPressure \\
\hline
ceiling & fixedValue $T_{\text{wall}}$ & noSlip & fixedFluxPressure \\
\hline
opposite\_wall & fixedValue $T_{\text{wall}}$ & noSlip & fixedFluxPressure \\
\hline
front, back & fixedValue $T_{\text{wall}}$ & noSlip & fixedFluxPressure \\
\hline
\end{tabularx}
}


\textbf{ヒーター熱流束の計算:}

$
q_{\text{heater}} = \frac{Q_{\text{kW}} \times 1000}{W_{\text{heater}} \times H_{\text{heater}}} \quad [\text{W/m}^2]
$

\subsubsection{4.2 セル単位の熱収支 (有限体積法)}

OpenFOAM は有限体積法 (FVM) で各セルのエネルギー収支を解く:

$
\underbrace{\frac{\partial}{\partial t} \int_{V_P} \rho h \, dV}_{\text{蓄熱}}
+ \underbrace{\sum_f (\rho \mathbf{U} h)_f \cdot \mathbf{S}_f}_{\text{対流フラックス}}
= \underbrace{\sum_f (\alpha_{\text{eff}} \nabla T)_f \cdot \mathbf{S}_f}_{\text{拡散フラックス}}
+ \underbrace{q \cdot V_P}_{\text{体積熱源}}
$

\begin{itemize} \item $f$: セル面, $\mathbf{S}_f$: 面積ベクトル \end{itemize}
\begin{itemize} \item $V_P$: セル体積 \end{itemize}
\begin{itemize} \item \textbf{蓄熱}: 非定常項（backward スキームで 2 次精度離散化） \end{itemize}
\begin{itemize} \item \textbf{対流フラックス}: 流れによる熱輸送（linearUpwind で離散化） \end{itemize}
\begin{itemize} \item \textbf{拡散フラックス}: 伝導+乱流拡散（Gauss linear corrected で離散化） \end{itemize}

\subsubsection{4.3 簡易版の熱収支模式図}

\begin{verbatim}
     ┌─────────────────────────────────┐
     │         上層 (T_upper)           │
     │                                 │
     │  入力: plume enthalpy           │
     │  出力: 壁面損失 (天井 + 上部側壁) │
     ├ ─ ─ ─ ─ ─ ─ z_int ─ ─ ─ ─ ─ ─ ┤  ← 界面 (質量保存で決定)
     │         下層 (T_lower)           │
     │                                 │
     │  入力: 輻射 + 界面伝導           │
     │  出力: 壁面損失 (床 + 下部側壁)   │
     │  出力: プルームへの空気供給       │
     │          ↑ プルーム              │
     │         [H] ヒーター              │
     └─────────────────────────────────┘
\end{verbatim}

\bigskip\hrule\bigskip

\subsection{5. 離散化スキームの対応表}

{\small
\noindent
\begin{tabularx}{\textwidth}{|l|L|L|L|}
\hline
\textbf{微分項} & \textbf{物理的意味} & \textbf{OpenFOAM スキーム} & \textbf{簡易版での扱い} \\
\hline
$\partial/\partial t$ & 時間変化 & backward (2次精度陰解法) & 前進 Euler ($\Delta t = 0.5$) \\
\hline
$\nabla \cdot (\rho \mathbf{U} \phi)$ & 対流 & linearUpwind (2次風上) & プルーム相関式で代替 \\
\hline
$\nabla \cdot (k \nabla T)$ & 拡散 & Gauss linear corrected & 界面伝導 ($k_{\text{int}}$) \\
\hline
$\nabla p$ & 圧力勾配 & Gauss linear & 解かない \\
\hline
$q_{\text{source}}$ & 熱源 & 境界条件 (面熱流束) & プルームエンタルピー \\
\hline
壁面熱伝達 & 壁損失 & 壁面関数 + fixedValue & Newton 冷却則 ($h_{\text{wall}}$) \\
\hline
\end{tabularx}
}


\bigskip\hrule\bigskip

\subsection{6. 数値解法と収束制御}

\subsubsection{6.1 正確版 — PIMPLE 法}

PIMPLE (PISO + SIMPLE の融合) アルゴリズム。
密閉空間の強い自然対流は本質的に非定常であるため、
定常ソルバー (SIMPLE) ではなく非定常ソルバーを採用し、
時間平均で統計量を取得する。

\begin{verbatim}
各時間ステップ:
  外側ループ (nOuterCorrectors = 2):
    1. 運動量方程式を解く → U* (仮の速度場)
    2. エネルギー方程式を解く → T
    3. 乱流方程式を解く → k, omega
    4. 密度を更新 → rho = f(p, T)
    内側ループ (nCorrectors = 1):
      5. 圧力補正方程式を解く → p'
      6. 速度を補正 → U = U* + U'
  時間ステップ調整 (maxCo = 0.5)
\end{verbatim}

\textbf{線形ソルバー:}

{\small
\noindent
\begin{tabularx}{\textwidth}{|l|L|L|L|}
\hline
\textbf{変数} & \textbf{ソルバー} & \textbf{前処理} & \textbf{許容残差} \\
\hline
$p_{rgh}$ & GAMG (代数マルチグリッド) & Gauss-Seidel & 1e-7 (rel 0.01) \\
\hline
$p_{rgh}$ Final & GAMG & Gauss-Seidel & 1e-7 (rel 0) \\
\hline
$U, T, k, \omega$ & PBiCGStab & DILU & 1e-7 (rel 0.1) \\
\hline
$U, T, k, \omega$ Final & PBiCGStab & DILU & 1e-7 (rel 0) \\
\hline
\end{tabularx}
}


\textbf{PIMPLE 残差制御:}

{\small
\noindent
\begin{tabularx}{\textwidth}{|L|L|}
\hline
\textbf{変数} & \textbf{閾値} \\
\hline
$p_{rgh}$ & 1e-4 \\
\hline
$U$ & 1e-4 \\
\hline
$T$ & 1e-5 (より厳密) \\
\hline
\end{tabularx}
}


\textbf{緩和係数 (Under-relaxation):}

{\small
\noindent
\begin{tabularx}{\textwidth}{|L|L|L|}
\hline
\textbf{変数} & \textbf{緩和係数} & \textbf{備考} \\
\hline
$p_{rgh}$ & 0.3 & 保守的 (圧力は不安定になりやすい) \\
\hline
$U$ & 0.7 & 中程度 \\
\hline
$T$ & 0.5 & やや保守的 (浮力との結合が強い) \\
\hline
$k, \omega$ & 0.7 & 中程度 \\
\hline
\end{tabularx}
}


\textbf{時間ステップ制御:}

{\small
\noindent
\begin{tabularx}{\textwidth}{|L|L|L|}
\hline
\textbf{パラメータ} & \textbf{値} & \textbf{備考} \\
\hline
deltaT & 0.05 s & 初期時間ステップ \\
\hline
maxCo & 0.5 & Courant 数上限 \\
\hline
adjustTimeStep & yes & 適応的時間ステップ調整 \\
\hline
endTime & 300 s & 計算終了時刻 \\
\hline
\end{tabularx}
}


\subsubsection{6.2 簡易版 — 擬似時間進行}

{\small
\noindent
\begin{tabularx}{\textwidth}{|L|L|}
\hline
\textbf{項目} & \textbf{値} \\
\hline
解法 & 前進 Euler（擬似時間ステップ $\Delta t = 0.5$） \\
\hline
状態変数 & $T_{\text{upper}}$, $T_{\text{lower}}$, $z_{\text{int}}$ \\
\hline
収束判定 & $\max(\lvert T_{\text{upper}}^{n+1} - T_{\text{upper}}^n \rvert, \; 10 \lvert z_{\text{int}}^{n+1} - z_{\text{int}}^n \rvert) < 10^{-4}$ \\
\hline
最大反復数 & 10,000 \\
\hline
物理クリッピング & $T_{\text{upper}} \in [T_{\text{wall}},\; T_{\text{wall}}+200]$ K \\
\hline
 & $z_{\text{int}} \in [0.05H,\; 0.95H]$ \\
\hline
\end{tabularx}
}


\bigskip\hrule\bigskip

\subsection{7. 物性値}

\subsubsection{空気 (全ソルバー共通)}

{\small
\noindent
\begin{tabularx}{\textwidth}{|l|L|L|L|}
\hline
\textbf{物性} & \textbf{記号} & \textbf{値} & \textbf{単位} \\
\hline
分子量 & $M_w$ & 28.96 & g/mol \\
\hline
定圧比熱 & $c_p$ & 1005 & J/(kg·K) \\
\hline
動粘性係数 & $\mu$ & $1.8 \times 10^{-5}$ & Pa·s \\
\hline
プラントル数 & $Pr$ & 0.7 & - \\
\hline
熱伝導率 & $\lambda = \mu c_p / Pr$ & 0.0258 & W/(m·K) \\
\hline
\end{tabularx}
}


\subsubsection{簡易版追加パラメータ}

{\small
\noindent
\begin{tabularx}{\textwidth}{|l|L|L|L|}
\hline
\textbf{パラメータ} & \textbf{記号} & \textbf{値} & \textbf{根拠} \\
\hline
基準密度 & $\rho_0$ & 1.1 kg/m³ & \textasciitilde{}300K での空気密度 \\
\hline
壁面熱伝達率 & $h_{\text{wall}}$ & 8.0 W/(m²·K) & 自然対流の標準値 \\
\hline
対流熱割合 & $f_{\text{conv}}$ & 0.7 & ヒーター出力の70\%が対流 \\
\hline
界面伝導率 & $k_{\text{int}}$ & 0.5 W/(m·K) & 界面を介した弱い伝導 \\
\hline
\end{tabularx}
}


\bigskip\hrule\bigskip

\subsection{8. 簡易版と正確版の対応関係}

\subsubsection{物理モデルの対応}

{\small
\noindent
\begin{tabularx}{\textwidth}{|l|L|L|L|}
\hline
\textbf{物理現象} & \textbf{正確版 (OpenFOAM)} & \textbf{簡易版 (2-Zone)} & \textbf{簡易化による影響} \\
\hline
\textbf{対流} & N-S 方程式で速度場を解く & MTT プルーム相関式で質量流量を算出 & 詳細な速度場は不明だが、プルームの総輸送量は再現 \\
\hline
\textbf{乱流} & SST k-omega で乱流粘性を計算 & プルーム相関式に乱流効果が内包 & 局所的な乱流強度の変化を再現不可 \\
\hline
\textbf{浮力} & 密度差 → 圧力勾配 → 速度場 & 密度の温度依存 + プルーム浮力相関 & 浮力-運動量の直接結合はないが、エネルギー輸送は再現 \\
\hline
\textbf{質量保存} & 連続の式で厳密に保存 & プルーム上昇 = 壁面沿い下降で保存 & ゾーン間の質量収支は保存される \\
\hline
\textbf{壁面境界層} & omegaWallFunction + メッシュ解像 & Newton 冷却 ($h_{\text{wall}}$ 固定) & 壁面近傍の詳細を無視 \\
\hline
\textbf{水平方向分布} & 3D で空間分解 & 各層内で均一を仮定 & 水平方向の温度むらを再現不可 \\
\hline
\textbf{密度変化} & 完全ガス $\rho = pM/RT$ & $\rho = \rho_0 T_{\text{wall}} / T$ & 簡易的だが温度依存あり \\
\hline
\end{tabularx}
}


\subsubsection{簡易版の限界}

\begin{enumerate} \item \textbf{速度場の不在}: プルーム相関式は総質量流量を与えるが、空間的な速度分布は得られない。 \end{enumerate}
   アウフグース（Phase 3）のジェット流れは原理的にモデル化不可能。
\begin{enumerate} \item \textbf{水平方向の均一仮定}: 各層内の水平温度分布が均一のため、 \end{enumerate}
   ヒーター直上と壁面近傍の温度差は再現できない。
\begin{enumerate} \item \textbf{蒸気輸送の困難}: ロウリュ（Phase 2）の蒸気は対流に乗って輸送されるが、 \end{enumerate}
   速度場がないため拡散モデルでの近似が必要になり精度に限界がある。

\subsubsection{簡易版で正しく再現できること}

\begin{itemize} \item \textbf{温度成層}: 上層高温・下層低温の 2 層構造（物理的に正しいメカニズムで再現） \end{itemize}
\begin{itemize} \item \textbf{界面高さ}: プルームの強さと壁面損失のバランスで決定される界面高さ \end{itemize}
\begin{itemize} \item \textbf{KPI の傾向}: K-01 (上下温度差) > 0 の判定は信頼できる \end{itemize}
\begin{itemize} \item \textbf{パラメータ感度}: ヒーター出力・壁温・部屋サイズを変えたときの応答傾向 \end{itemize}
\begin{itemize} \item \textbf{エントレインメント効果}: プルーム高さが大きいほど巻き込みが増え温度が低下する現象 \end{itemize}

\bigskip\hrule\bigskip

\subsection{9. 既知の物理的制約とフェーズ拡張計画}

本節では、初期 Phase 1 の方程式系に内在していた物理的制約と、
その後のフェーズで実施した方程式拡張を記録する。
項目 9.1--9.4 は全て実装済であり、現在の方程式系（§3.3, §10, 付録C）に反映されている。

\subsubsection{9.1 蒸気化による体積膨張 \textnormal{【解決済: Phase 2 で実装】}}

\textbf{旧問題:} 界面質量保存式が乾燥空気のプルームのみを扱い、
ロウリュ時の蒸気体積膨張（$\dot{V}_{\text{steam}}$）が欠落していた。

\textbf{実装内容:} 過渡ソルバー (\texttt{solve\_transient}) に以下を実装:

$
\dot{V}_{\text{steam}} = \frac{\dot{m}_{\text{water}} \, R \, T_{\text{upper}}}{p \, M_{w,\text{H}_2\text{O}}}
$

\begin{itemize} \item $\dot{m}_{\text{water}}$: Spalding 型指数減衰蒸発モデル ($\dot{m} = m_0/\tau \cdot e^{-t/\tau}$) から区間積分で算出 \end{itemize}
\begin{itemize} \item 蒸発潜熱はヒーター石側から奪われるため、空気への正味エネルギー追加はない \end{itemize}
\begin{itemize} \item 界面質量保存式（§3.3）に $\dot{V}_{\text{steam}}$ 項として反映済 \end{itemize}

\subsubsection{9.2 アウフグース ROM + 質量保存 \textnormal{【解決済: Phase 3 で実装】}}

\textbf{旧問題:} エネルギー交換のみモデル化し、質量保存式への反映が欠落していた。

\textbf{実装内容:} ROM 強制混合 ($\beta_{\text{aug}}$ [kg/s]) を
エネルギー保存と界面質量保存の両方に実装。

\paragraph{エネルギー交換}
上層: $\Delta T_{\text{upper}} \mathrel{-}= \beta_{\text{aug}} \, c_p \, (T_{\text{upper}} - T_{\text{lower}}) / (m_{\text{upper}} \, c_p)$

下層: $\Delta T_{\text{lower}} \mathrel{+}= \beta_{\text{aug}} \, c_p \, (T_{\text{upper}} - T_{\text{lower}}) / (m_{\text{lower}} \, c_p)$

\paragraph{質量保存（界面方程式への寄与）}
双方向質量交換の密度差による正味体積効果:

$
\dot{V}_{\text{mix}} = \beta_{\text{aug}} \left( \frac{1}{\rho_{\text{lower}}} - \frac{1}{\rho_{\text{upper}}} \right)
$

$\rho_{\text{upper}} < \rho_{\text{lower}}$ のとき $\dot{V}_{\text{mix}} < 0$ → 界面降下
（タオルが高温空気を強制的に下方へ押す効果）。
界面質量保存式（§3.3）に反映済。

\subsubsection{9.3 正確版の多成分化: 組成浮力の導入 \textnormal{【未実装: Phase 2+ 予定】}}

\textbf{問題:} 現在の \texttt{pureMixture} + \texttt{perfectGas}（$M_w = 28.96$ 固定）は、
水蒸気（$M_w = 18.015$）と乾燥空気の混合による密度変化を表現できない。
ロウリュ時のプルーム急上昇は、温度差による \textbf{熱的浮力} に加え、
蒸気濃度差による \textbf{組成浮力（solutal buoyancy）} の寄与が大きい。
水蒸気は空気より軽い（$18/29 \approx 0.62$）ため、蒸気リッチなプルームは
温度差だけで予測されるより速く上昇する。

\textbf{拡張方針:} Phase 2 で熱物理モデルを変更:

\begin{verbatim}
thermoType
{
    type            heRhoThermo;
    mixture         multiComponentMixture;   // ← pureMixture から変更
    transport       sutherland;              // ← 温度依存粘性
    thermo          janaf;                   // ← 温度依存 Cp
    equationOfState perfectGas;
    specie          specie;
    energy          sensibleEnthalpy;
}
\end{verbatim}

蒸気質量分率 $Y$ の輸送方程式を追加:

$
\frac{\partial (\rho Y)}{\partial t} + \nabla \cdot (\rho \mathbf{U} Y)
= \nabla \cdot (D_{\text{eff}} \nabla Y) + S_Y
$

混合密度:

$
\rho = \frac{p}{R T} \left( \frac{Y}{M_{w,\text{H}_2\text{O}}} + \frac{1-Y}{M_{w,\text{air}}} \right)^{-1}
$

ここで $D_{\text{eff}} = D_{\text{mol}} + D_t$（分子拡散 + 乱流拡散）。
$D_{\text{mol}} \approx 2.5 \times 10^{-5}$ m²/s（空気中の水蒸気拡散係数, 373K）。

これにより密度が温度 $T$ と蒸気質量分率 $Y$ の両方に依存し、
組成浮力が自動的に運動量方程式に反映される。

\subsubsection{9.4 輻射モデルの改善 \textnormal{【簡易版: 解決済 / OpenFOAM: 未実装】}}

\textbf{旧問題:} 簡易版の下層エネルギー保存における $Q_{\text{rad}} \times 0.3$（固定係数）は
以下の感度を失っていた:

\begin{itemize} \item ストーブの配置（部屋の中央 vs 隅） \end{itemize}
\begin{itemize} \item 部屋の形状・アスペクト比 \end{itemize}
\begin{itemize} \item ベンチや障害物による遮蔽 \end{itemize}

サウナにおける人体の受熱量の半分以上は輻射によるものであり（特に高温域）、
固定係数では熱ストレス指標（K-06）の空間的分布を正しく予測できない。

\textbf{簡易版の実装内容:}

\texttt{\_compute\_view\_factors()} で立体角近似による形態係数を算出。
ヒーターから床、下部壁、上部壁、天井、人体への5分割で閉合則を満たす。

$
Q_{\text{rad,lower}} = Q_{\text{rad}} \times (F_{\text{heater} \to \text{floor}} + F_{\text{heater} \to \text{lower walls}})
$

さらに、ヒーターから人体への直接放射 $q_{\text{rad,body}}$ を
小面積近似 ($F_{\text{body}} \approx A_{\text{body}} / (2\pi d^2)$) で算出し、
皮膚熱収支モデル（付録C）に反映。部屋の形状やヒーター配置の変更に対して
自動的に輻射配分が変化する。

\textbf{拡張方針 (OpenFOAM):}

Phase 2 以降で輻射モデルを導入する:

\begin{verbatim}
radiationModel  fvDOM;    // Finite Volume Discrete Ordinates Method
// または
radiationModel  viewFactor;  // 壁面間のみの輻射交換（計算コスト低）

absorptionEmissionModel  greyMeanAbsorptionEmission;
scatterModel    none;
\end{verbatim}

特に高温域（$T > 100$ °C）では壁面輻射による二次加熱（ストーブ → 壁 → 空気層 → 反対側の壁）
が自然対流の境界層形成に影響するため、Phase 2 では必須。
\texttt{viewFactor} モデルは壁面間のみの輻射交換を計算するため、計算コストが低く
PoC フェーズに適している。

\bigskip\hrule\bigskip

\subsection{10. 換気モデル (Natural Ventilation)}

\subsubsection{10.1 モデル概要}

実サウナには給気口（床付近・ヒーター近傍）と排気口（天井付近）があり、
室内外温度差による煙突効果（stack effect）で自然換気が発生する。
従来の密閉箱仮定（sealed-box assumption）を除去し、
任意の給排気構成をサポートする。

\begin{verbatim}
 ┌──────────────────────────────────┐  ← ceiling
 │         上層 (T_upper)            │
 │                            [排気] ← z_exhaust (排気口高さ)
 ├ ─ ─ ─ ─ ─ ─ z_int ─ ─ ─ ─ ─ ─ ─┤  ← 界面
 │         下層 (T_lower)            │
 │  [給気] ←  z_supply (給気口高さ)   │
 │          ↑ プルーム               │
 │         [H] ヒーター               │
 └──────────────────────────────────┘  ← floor
     T_ambient (外気温)
\end{verbatim}

\texttt{model: "none"} (デフォルト) で従来の密閉モデル、
\texttt{model: "natural"} で自然換気モデルを選択。

\subsubsection{10.2 煙突効果による駆動圧力}

給気口と排気口の高さ差による温度差駆動圧力（stack pressure）:

$
\Delta p = \rho_{\text{amb}} \, g \, (z_{\text{exhaust}} - z_{\text{supply}}) \,
\frac{\bar{T}_{\text{col}} - T_{\text{ambient}}}{T_{\text{ambient}}}
$

ここで $\bar{T}_{\text{col}}$ は給排気口間の高さ加重平均温度:

$
\bar{T}_{\text{col}} = f_{\text{lower}} \, T_{\text{supply}} + f_{\text{upper}} \, T_{\text{exhaust}}
$

$f_{\text{lower}} = \frac{\min(z_{\text{int}}, z_{\text{exhaust}}) - z_{\text{supply}}}{z_{\text{exhaust}} - z_{\text{supply}}}$, \quad
$f_{\text{upper}} = 1 - f_{\text{lower}}$

\begin{itemize} \item $\rho_{\text{amb}}$: 外気密度 [kg/m³] \end{itemize}
\begin{itemize} \item $z_{\text{exhaust}}, z_{\text{supply}}$: 排気口・給気口の高さ [m] \end{itemize}
\begin{itemize} \item $T_{\text{supply}}, T_{\text{exhaust}}$: 各口位置のゾーン温度（界面位置に応じて $T_{\text{upper}}$ または $T_{\text{lower}}$） \end{itemize}
\begin{itemize} \item $T_{\text{ambient}}$: 外気温 [K] \end{itemize}

\subsubsection{10.3 オリフィス流量モデル}

有効開口面積を用いた質量流量:

$
\dot{m}_{\text{vent}} = C_d A_{\text{eff}} \sqrt{2 \rho_{\text{supply}} |\Delta p|}
$

ここで:

\begin{itemize} \item $C_d A_{\text{eff}} = \min(C_{d,s} A_s, \, C_{d,e} A_e)$: 有効開口面積（流量は狭い方で制限） \end{itemize}
\begin{itemize} \item $\rho_{\text{supply}}$: 給気口位置の空気密度 [kg/m³] \end{itemize}
\begin{itemize} \item $\dot{m}_{\text{vent}}$: 正（流入）が通常だが、ロウリュ蒸気膨張時の室内圧上昇で負（逆流＝排出）も許容 \end{itemize}

\subsubsection{10.4 エネルギー保存への寄与}

\paragraph{下層（給気側）}

外気が $T_{\text{ambient}}$ で流入し、下層空気と混合:

$
Q_{\text{vent,lower}} = \dot{m}_{\text{vent}} \, c_p \, (T_{\text{ambient}} - T_{\text{lower}})
$

（$T_{\text{ambient}} < T_{\text{lower}}$ のとき冷却効果）

\paragraph{上層（排気側）}

排気口から $T_{\text{upper}}$ の空気が排出され、界面から $T_{\text{lower}}$ の空気で置換:

$
Q_{\text{vent,upper}} = \dot{m}_{\text{vent}} \, c_p \, (T_{\text{lower}} - T_{\text{upper}})
$

（常に冷却効果: 高温の上層空気が低温の下層空気で置換される）

\subsubsection{10.5 界面高さへの寄与}

給気は下層に体積を追加し、界面を上方に押し上げる:

$
A_{\text{floor}} \, \frac{dz_{\text{int}}}{dt}
= -\frac{\dot{m}_p}{\rho_{\text{upper}}} + \dot{V}_{\text{return}} - \dot{V}_{\text{steam}}
+ \frac{\dot{m}_{\text{vent}}}{\rho_{\text{lower}}}
$

定常状態では、プルームによる界面降下と換気・壁面還流による界面上昇がバランスする。

\subsubsection{10.6 湿度への寄与（過渡計算時）}

外気は低湿度（$w_{\text{ambient}}$）であり、ロウリュ後の湿度減衰メカニズムとなる:

$
\frac{dw}{dt} = \frac{\dot{m}_{\text{vent}}}{m_{\text{upper}}} (w_{\text{ambient}} - w)
$

この1次減衰により、換気が活発なほどロウリュ後の湿度ピークからの回復が早くなる。

\subsubsection{10.7 YAML 構成例}

\begin{verbatim}
ventilation:
  model: "natural"       # "natural" or "none" (default)
  supply:
    height: 0.3          # m above floor
    area: 0.02           # m^2 (e.g. 20cm x 10cm vent)
    Cd: 0.6              # discharge coefficient
  exhaust:
    height: 2.0          # m above floor (near ceiling)
    area: 0.02           # m^2
    Cd: 0.6
  T_ambient: 293.15      # K (outside temperature)
  w_ambient: 0.005       # kg/kg (outside humidity ratio)
\end{verbatim}

\bigskip\hrule\bigskip

\subsection{付録 A: フェーズ別の実装状況}

{\small
\noindent
\begin{tabularx}{\textwidth}{|l|L|L|L|}
\hline
\textbf{フェーズ} & \textbf{OpenFOAM} & \textbf{簡易版} & \textbf{状態} \\
\hline
Phase 1 & buoyantPimpleFoam + SST k-omega ($G_b$ 浮力生成項付), pureMixture & 2-Zone プルームモデル (定常) + Heskestad 仮想原点補正 & \textbf{実装済} \\
\hline
Phase 2 & multiComponentMixture + 蒸気輸送 ($Y$) + H2O テンプレート & 蒸気体積膨張 $\dot{V}_{\text{steam}}$ + 形態係数輻射 + 壁面昇温モデル + 湿度連成物性 + 皮膚熱収支モデル + 自然換気モデル + 非定常ソルバー (\texttt{solve\_transient}) & \textbf{実装済} \\
\hline
Phase 3 & 運動量ソース項 (ジェット) — 未実装 & ROM: 強制混合係数 $\beta_{\text{aug}}$ + 質量保存整合 ($\dot{V}_{\text{mix}}$) + 抽出スクリプト (\texttt{extract\_beta\_aug.py}) & \textbf{枠組み実装済} \\
\hline
Phase 4-5 & — & CSV 取込 (\texttt{validation.py}) + プローブ比較 (\texttt{compare\_probes}) + レポート生成 (\texttt{reporting.py}) & \textbf{枠組み実装済} \\
\hline
\end{tabularx}
}


\subsection{付録 B: KPI 一覧と実装状況}

{\small
\noindent
\begin{tabularx}{\textwidth}{|l|L|L|L|}
\hline
\textbf{KPI ID} & \textbf{名称} & \textbf{入力} & \textbf{実装} \\
\hline
K-01 & 定常温度差 (上段-下段) & プローブ定常値 & \textbf{済} (\texttt{kpi.py}) \\
\hline
K-02 & Löyly 後ピーク温度上昇 & 時系列 $T_{\text{upper}}(t)$ & \textbf{済} \\
\hline
K-03 & Löyly 後ピーク絶対湿度 & 時系列 $w(t)$ & \textbf{済} \\
\hline
K-04 & ピーク到達時間 & 時系列 $T_{\text{upper}}(t)$ + イベント時刻 & \textbf{済} \\
\hline
K-05 & 顔面風速ピーク (proxy) & $\beta_{\text{aug}}$ から推定 & \textbf{済} (ROM proxy) \\
\hline
K-06 & 簡易熱ストレス指標 & 体感温度 (皮膚熱収支モデル+放射) & \textbf{済} \\
\hline
K-07 & 上下相対温度差 & プローブ定常値 & \textbf{済} (\texttt{kpi.py}) \\
\hline
\end{tabularx}
}


\subsection{付録 C: 簡易版の追加モデル (Phase 2 実装分)}

\subsubsection{壁面集中定数モデル (lumped wall)}

壁面内面温度 $T_{\text{wall,inner}}$ を状態変数として追加:

$
(\rho c_p)_w \, \delta_w \, A_{\text{wall}} \, \frac{dT_{\text{wall,inner}}}{dt}
= \underbrace{Q_{\text{conv→wall}}}_{\text{空気→壁面対流}}
+ \underbrace{Q_{\text{rad→wall}}}_{\text{ヒーター輻射}}
- \underbrace{\frac{\lambda_w}{\delta_w} A_{\text{wall}} (T_{\text{wall,inner}} - T_{\text{wall,outer}})}_{\text{壁面→外部熱伝導}}
$

パラメータ:

\begin{itemize} \item $\delta_w = 0.015$ m (木製パネル厚さ) \end{itemize}
\begin{itemize} \item $\lambda_w = 0.12$ W/(m·K) (木の熱伝導率) \end{itemize}
\begin{itemize} \item $(\rho c_p)_w = 5 \times 10^5$ J/(m³·K) (木の体積熱容量) \end{itemize}

\texttt{model: fixed} で従来の固定壁温、\texttt{model: lumped} で昇温モデルを選択。

\subsubsection{湿度連成物性モデル}

混合気体の物性を湿度比 $w$ [kg/kg] に応じて補正:

$
c_{p,\text{mix}} = (1 - y_v) \, c_{p,\text{air}} + y_v \, c_{p,\text{vapor}}
$

$
h_{\text{wall,eff}} = h_{\text{wall,base}} \left(\frac{c_{p,\text{mix}}}{c_{p,\text{air}}}\right)^{0.25} \left(\frac{\lambda_{\text{mix}}}{\lambda_{\text{air}}}\right)^{0.75}
$

ここで $y_v = w / (1+w)$ は蒸気質量分率。

\subsubsection{体感温度 (皮膚熱収支モデル)}

従来の Steadman (1979) 近似は屋外条件 ($20$--$50$\textdegree C) のみ有効であり、
サウナ温度域 ($60$--$120$\textdegree C) では適用外となる。
本モデルでは皮膚表面の熱収支を直接評価し、等価温度 $T_{\text{eq}}$ を算出する。

\paragraph{対流熱流束}
$
q_{\text{conv}} = h_{\text{conv}} \, (T_{\text{air}} - T_{\text{skin}})
$

ここで $T_{\text{skin}} = 36$\textdegree C (平均皮膚温度)、
$h_{\text{conv}} = 8 \; \text{W/(m}^2\text{K)}$ (自然対流熱伝達率)。

\paragraph{蒸発 / 凝縮項}
Lewis 関係を用い蒸発冷却 (または高湿時の凝縮加熱) を算出する。
皮膚濡れ率 $W = 0.4$（サウナ環境の典型値）で蒸発能力を制限し、
生理学的上限 $q_{\text{evap,max}} = 400$ W/m² を設ける。

$p_{\text{vapor}} > p_{\text{sat,skin}}$ (凝縮) のとき（$W$ は不適用）:
$
q_{\text{evap}} = \frac{16.5 \, h_{\text{conv}} \, (p_{\text{vapor}} - p_{\text{sat,skin}})}{1000}
$

$p_{\text{vapor}} \leq p_{\text{sat,skin}}$ (蒸発冷却) のとき:
$
q_{\text{evap,raw}} = \frac{16.5 \, h_{\text{conv}} \, (p_{\text{sat,skin}} - p_{\text{vapor}})}{1000}
$
$
q_{\text{evap}} = -\min(W \cdot q_{\text{evap,raw}}, \; q_{\text{evap,max}})
$

ここで $p_{\text{sat}}$ は Magnus 式から算出 ($p_{\text{sat}} = 610.78 \exp(17.27 T / (T + 237.3))$ [Pa])。

\paragraph{ヒーターからの直接放射}
ヒーター表面温度を放射出力と面積から推定し、ヒーター--人体間の形態係数 $F_{\text{body}}$
を用いて人体への放射熱流束を算出する。

ヒーター表面温度の推定:

$
T_{\text{heater}} = \left( \frac{q_{\text{heater,surface}}}{\varepsilon_{\text{heater}} \, \sigma} + T_{\text{wall,inner}}^4 \right)^{1/4}
$

ここで $q_{\text{heater,surface}} = Q_{\text{rad}} / A_{\text{heater}}$ [W/m²]、
$\varepsilon_{\text{heater}} = 0.90$（ストーブ/金属表面放射率）。

人体への正味放射熱流束:

$
q_{\text{rad,body}} = \varepsilon_{\text{body}} \, \sigma \, F_{\text{body}} \, (T_{\text{heater}}^4 - T_{\text{skin}}^4)
$

ここで $\varepsilon_{\text{body}} = 0.97$（人体皮膚放射率）、
$\sigma = 5.67 \times 10^{-8} \; \text{W/(m}^2\text{K}^4\text{)}$。

形態係数 $F_{\text{body}}$ はヒーター中心から上段ベンチ上の人体までの距離を用い、
小面積近似 ($F \approx A_{\text{body}} / (2\pi d^2)$) で算出する。
$F_{\text{body}}$ は [0.005, 0.10] にクリップされる。

\paragraph{等価温度}
全熱流束の合計を基準熱伝達率 $h_{\text{ref}} = 10 \; \text{W/(m}^2\text{K)}$ で正規化する。

$
T_{\text{eq}} = T_{\text{skin}} + \frac{q_{\text{conv}} + q_{\text{rad,body}} + q_{\text{evap}}}{h_{\text{ref}}}
$

この等価温度は K-06 (簡易熱ストレス指標) に直接使用される。
$60$--$120$\textdegree C の範囲で物理的に妥当な値を返す。

\subsubsection{ヒーター容量の目安}

業界標準 (Harvia, HELO): 約 \textbf{1 kW / m³} のサウナ室体積。

{\small
\noindent
\begin{tabularx}{\textwidth}{|L|L|L|}
\hline
\textbf{部屋サイズ} & \textbf{体積} & \textbf{推奨ヒーター} \\
\hline
2.0×1.5×2.2 m (1-2人) & 6.6 m³ & 6-8 kW \\
\hline
3.0×2.5×2.5 m (4-6人) & 18.8 m³ & 18-20 kW \\
\hline
4.0×3.0×2.5 m (商業用) & 30 m³ & 30 kW \\
\hline
\end{tabularx}
}



\end{document}
